\section{Accuracy rule or cost rule? }
\label{sec:sweep}
\textbf{Code:} \code{experiments/task2\_cost\_sweep.py}.

With both classifiers of Section~\ref{sec:models} in place, we applied the rule of
Section~\ref{sec:rule} to them and asked how much it beats the accuracy rule as the two error costs
move apart. The sweep varies the ratio $r=\cfp/\cfn$ from $1/20$ to $20$ (with
$\cfp=\sqrt r$, $\cfn=1/\sqrt r$ and $\ci=0.1$) and applies three rules to the same fitted
probabilities: the accuracy rule (spam iff $\post\ge0.5$), the cost rule without inspection (spam iff
$\post\ge\tstar$) and the full rule with inspection. It runs logistic regression and both Gaussian
models on Spambase (standardised, no $\log(1+x)$) and on a synthetic five-dimensional data set whose
two classes have different covariances. On the synthetic data an \emph{oracle} applies the same rules
to the exact posterior.

\begin{figure}[!htbp]
\centering
\includegraphics[width=\linewidth]{task2_loss_spambase.png}
\caption{Spambase: realized test loss of the three rules against $\cfp/\cfn$ (log scale) for each
model, mean $\pm$ standard error over 20 splits, $\ci=0.1$.}
\label{fig:t2}
\end{figure}

\begin{table}[!htbp]
\centering\small
\input{\resultsdir/task2_loss_spambase.tex}
\caption{Spambase: realized test loss at three cost ratios (mean $\pm$ standard error, 20 splits).}
\label{tab:t2}
\end{table}

\begin{table}[!htbp]
\centering\small
\input{\resultsdir/task2_loss_synthetic.tex}
\caption{Synthetic data: realized test loss at three cost ratios. ``Oracle'' applies the rule with
inspection to the exact posterior.}
\label{tab:t2-syn}
\end{table}

Figure~\ref{fig:t2} and Tables~\ref{tab:t2} and~\ref{tab:t2-syn} show three things.
\begin{itemize}
  \item \emph{The accuracy rule is right only at equal costs.} At $r=1$ it makes the same decisions as
        the cost rule. At $r=0.1$ it costs logistic regression $0.166$ per email against $0.082$, and
        at $r=10$ $0.115$ against $0.082$.
  \item \emph{Inspection is the cheapest rule in every column}, for every model. For logistic
        regression it brings the loss to $0.050$--$0.053$ while inspecting $23$--$32\%$ of emails.
  \item \emph{Logistic regression is the best model on Spambase} with every rule at all three
        ratios. With the two cost-based rules the class-specific Gaussian is the worst. On the synthetic
        data, where the class-specific Gaussian is the correct model, it is within $0.001$ of the oracle
        at all three ratios (for example $0.077$ against $0.076$ at $r=0.1$), while logistic regression,
        which can only draw a linear boundary, is up to $0.023$ worse.
\end{itemize}

To see which emails go to the inspector, Table~\ref{tab:t2-conf} breaks the logistic model's
decisions down by true label (counts per test split of about $920$ emails). As real emails become
dearer to lose ($r=10$), the rule without inspection files fewer of them as spam ($9.7$ against $28.8$
for the accuracy rule) but has to deliver $141.6$ spam to do so. With inspection it files only $6.3$
real emails as spam and delivers $24.4$ spam, by sending $209$ uncertain emails, mostly spam, to the
inspector. Figure~\ref{fig:t2-rej} shows how the share of inspected emails changes with the ratio.

\begin{table}[!htbp]
\centering\footnotesize
\input{\resultsdir/task2_confusion_spambase_logistic.tex}
\caption{Spambase, logistic regression: mean confusion matrices (true label $\times$ action) per test
split, $\ci=0.1$.}
\label{tab:t2-conf}
\end{table}

\begin{figure}[!htbp]
\centering
\includegraphics[width=\linewidth]{task2_reject_spambase.png}
\caption{Spambase: share of emails inspected by the rule with inspection against $\cfp/\cfn$, for each
model ($\ci=0.1$).}
\label{fig:t2-rej}
\end{figure}

The rule compares $\post$ with thresholds far from $0.5$, so a model's cost depends on how far its
probabilities can be trusted, not only on which side of $0.5$ they fall. That is what we checked next.
